\chapter{Conclusions}\label{demo-outlook}
The thesis frames sensor placement as a transparent comparison among limited choices. Its contribution is an original worked narrative and inspectable synthetic document, not a validated monitoring system.
\section{Synthesis}\label{demo-outlook-summary}
Coverage, warning delay, and resilience answer different questions. Keeping them separate makes the consequences of a placement rule easier to discuss and avoids treating a downstream observation as equivalent to an early upstream warning.
\subsection{Conclusions within the model}\label{demo-outlook-answer}
Under the imagined branching topology, spreading sensors can protect distinct source paths. A shared downstream sensor can provide backup. Neither strategy dominates without specifying the event schedule, observation window, and outage assumptions.
\section{From illustrative networks to a calibrated environmental monitoring study}\label{demo-outlook-next}
A genuine study would first implement the specified event accounting and verify it against hand-worked networks. Only then would expanded topology families and uncertain travel times provide meaningful additional comparisons.
\subsection{Independent transport calibration}\label{demo-outlook-calibration}
Connecting the model to a real catchment would require independent measurements of transport and mixing. Choosing a parameter because it improves a placement result would confound calibration with the desired conclusion.
\subsubsection{Separating calibration from placement comparisons}\label{demo-outlook-separation}
The data used to choose model parameters should be distinguished from the scenarios used to compare placement rules. This separation makes it easier to identify which conclusion depends on which assumption.
\begin{figure}[htbp]
\centering
\fbox{\begin{minipage}{0.8\linewidth}
Inspect miniature cases $\longrightarrow$ implement event accounting\\[4pt]
Challenge assumptions $\longrightarrow$ gather independent evidence\\[4pt]
Evaluate feasibility $\longrightarrow$ consider a deployment decision
\end{minipage}}
\caption{A conditional path beyond this fictional document; none of these later activities is claimed as completed research.}\label{demo-fig-outlook}
\end{figure}
\section{Closing remarks}\label{demo-outlook-closing}
Listening to an imaginary river is useful when it clarifies what a real experiment would need to establish. The boundary between an inspectable example and a supported claim is part of the research design, not an afterthought.
